
\documentclass[12pt]{article}
\usepackage[a4paper, margin=1in]{geometry}
\usepackage{graphicx} 
\usepackage{amsmath} 
\usepackage{amssymb} 
\usepackage{dsfont}
\usepackage{comment}
\begin{document}
\section{Notation}
We model four events each year: growth, colonization, fizzle and burnout. \\

Let $(k,t)$ be the time index, where $k$ represent the year index and $t$ represent the season index, where $t=0$ for the time at the beginning of the year, $t=1$ for the time after growth, $t=2$ for the time after colonization, $t=3$ for the time after fizzle and $t=4$ for the time after burnout. So $(k,0)$ represents the same time as $(k-1,4)$\\

Assume there are $n$ patches and we label them with $p_i$. For each patch we use $(\mathds{1}_i^{(k,t)},N_i^{(k,t)})$ to describe its state. $N_i^{(k,t)}$ represents population size in patch $p_i$ at time point $(k,t)$, and $\mathds{1}_i^{(k,t)}$ is defined as:
$$
\mathds{1}_i^{(k,t)} =
\begin{cases}
1, & \text{if $p_i$ is infected at time $(k,t)$}, \\
0, & \text{otherwise} 
\end{cases}
$$

Let $U \sim \text{Uniform}(0, 1)$ be a random variable uniformly distributed over $(0,1)$ , and we have
$$
\mathds{1}_{U \leq P} =
\begin{cases}
1, & \text{if } U \leq P, \\
0, & \text{otherwise.}
\end{cases}
$$

\section{Growth}
We can assume that reproduction of rats follows logistic growth, i.e.,
$$N_i^{(k,1)}=N_i^{(k,0)}+rN_i^{(k,0)}(1-\frac{N_i^{(k,0)}}{K})$$
where $r$ is the growth rate and $K$ is the carry capacity of a patch. We currently assume that all patches are of equal area and have the same carrying capacity.

\section{Colonization}
\subsection{Population size}
We assume that each individual rat move with probability $c$ and land in a random patch. We ignore different distances and assume that all patches are of equal availability to a moving rat. So 

$$N_i^{(k,2)}=(1-c)N_i^{(k,1)}+\frac{c\sum _iN_i^{(k,1)}}{n}$$
The first term is population staying in the patch and the second term is population moving into the patch. 

\subsection{Infection}
\subsubsection{version 1 (based on death number)}
When an infected host dies the fleas are released and can travel between patches (such as being carried by a moving individual) and infect hosts in other patches. \\

Let $z$ be the final size. The number of individuals that died of the disease in the previous year in a patch is $zDN_i^{(k-1,3)}$ if patch $p_i$ is infected, and  0 if not. So the total number of death is $\sum_i zDN_i^{(k-1,3)}\mathds{1}_i^{(k-1,3)}$. \\

Assume the probability of a susceptible patch getting infected follows hazard model  
$$P_{\rm{I}}^{(k)}=1-{\mathrm{e}}^{-\frac{\alpha zD\sum_iN_i^{(k-1,3)}\mathds{1}_i^{(k-1,3)}}{nK}}$$
where $\alpha$ is a scaling constant.

\subsubsection{version 2 (based on colonization)}
Let $z$ be the final size. The number of individuals that was infected in the previous year in a patch is $zN_i^{(k-1,3)}$ if patch $p_i$ is infected, and  0 if not, and they move with probability $c$. So the total number of infected rats moving is 
$$cz\sum_i N_i^{(k-1,3)}\mathds{1}_i^{(k-1,3)}$$
and the expected number of infected rats moving into a patch is 

$$\frac{cz\sum_i N_i^{(k-1,3)}\mathds{1}_i^{(k-1,3)}}{n}$$

Assume the probability of a susceptible patch getting infected follows hazard model  
$$P_{\rm{I}}^{(k)}=1-{\mathrm{e}}^{-\frac{\alpha cz\sum_i N_i^{(k-1,3)}\mathds{1}_i^{(k-1,3)}}{n}}$$
where $\alpha$ is the probability of a susceptible patch getting infected caused by an infected individual moving in.\\

And
$$\mathds{1}_i^{(k,2)}=1-(1-\mathds{1}_i^{(k,1)})\mathds{1}_{U<P_{\rm{I}}^{(k)}}$$

\section{Fizzle}
Fizzle probability is 
$$P_{\rm{f}}=\frac{1}{\mathcal{R}_0}$$
And
$$ \mathds{1}_i^{(k,3)}= \mathds{1}_i^{(k,2)} \mathds{1}_{U<P_{\rm{f}}}$$
Population size does not change.

\section{Burnout}
In infected patches, the population should be reduced:

$$ N_i^{(k,4)}=(1-zD\mathds{1}_i^{(k,3)})N_i^{(k,3)} $$
Burnout probability if a patch didn't fizzle is 
$$P_{\rm{b}}=q^{Ny^\star}$$
($q$ and $y^\star$ are defined in burnout paper)

And
$$ \mathds{1}_i^{(k,4)}= \mathds{1}_i^{(k,3)} \mathds{1}_{U<P_{\rm{b}}}$$

\end{document}

